% !TEX root = ../thesis.tex

Legged robots are increasingly deployed in unstructured environments where
the ability to \emph{manipulate} objects is as important as the ability to
\emph{locomote}.
This project addresses the challenge of whole-body loco-manipulation by
reproducing and adapting the RoboDuet cooperative control framework
(Pan et al., IEEE RA-L 2025) to the Unitree B2 quadruped equipped with a
Unitree Z1 robotic arm --- a platform four times heavier than the Go1
for which the original framework was designed.

We identify a critical PD-gain mismatch as the principal obstacle to
direct framework transfer and resolve it by increasing the leg joint
stiffness from $K_p = 35$ to $K_p = 200$\,N\,m/rad and the damping
from $K_d = 1$ to $K_d = 20$\,N\,m\,s/rad.
Building on this corrected baseline, we implement a three-stage
reinforcement-learning curriculum using Proximal Policy Optimisation
in the Isaac Gym simulator with 2048 parallel environments:
(1)~a locomotion pre-training stage in which the dog policy learns
stable trotting on flat terrain with the arm fixed;
(2)~an arm pre-training stage in which the arm policy learns 6-DoF
end-effector pose tracking with the body standing still; and
(3)~a hybrid fine-tuning stage that hot-starts from both pre-trained
policies and trains the cooperative walk-while-grasp behaviour.

The cooperative architecture employs separate actor-critic networks for
the locomotion and arm policies.
The arm policy outputs both six arm joint targets and a body
orientation command (pitch, roll) that is injected into the locomotion
policy's observation vector, enabling the body to lean dynamically and
expand the arm's reachable workspace.
Isaac Gym simulation results demonstrate that the B2+Z1 system can
sustain stable locomotion at commanded velocities while simultaneously
tracking arm end-effector targets, validating the adapted framework
on the heavier platform.
